

\footnotesize%
\pstart%
\noindent%
Das Stück N.\;\ref{60270} ist auf Papier geschrieben, wie es Leibniz für die Abfassung der \textit{Dynamica}\cite{01354} während seines Aufenthaltes in Rom\indexO{Rom} verwendete. Dort traf er am 14. April 1689 ein, so dass N.\;\ref{60270} erst danach entstanden sein kann.  Darin finden sich zwei Verweise auf die \textit{Dynamica}\cite{01354}, die Leibniz beide nachträglich ergänzt und präzisiert, einmal unter Angabe der Proposition mit genauer Nummer, das andere Mal mit Bezug auf eine nicht gezählte Proposition im zweiten Teil der \textit{Dynamica}\cite{01354} (\refpassage{LH_37_04_169-170_21a}{LH_37_04_169-170_21b} und \refpassage{LH_37_04_169-170_22a}{LH_37_04_169-170_22b}). 
\pend
\pstart
Noch auf der Rückreise nach Deutschland schreibt Leibniz am 20. Februar 1690 aus Venedig\indexO{Venedig} an Rudolf Christian von Bodenhausen\indexN{\textls{Bodenhausen} (Bodenus, Bodenausen), Rudolf Christian, Freiherr v. 1698}: Er kündigt ihm weitere Teile für die \textit{Dynamica}\cite{01354} an, um dieses Werk für eine Reinschrift zu vervollständigen, darunter auch eine Schrift über den „Modum construendi Thermometrum accurate divisum“. Dabei handelt es sich aller Wahrscheinlichkeit nach um das vorliegende Stück N.\;\ref{60270}. Das „päckchen“ mit diesem und weiteren Nachträgen zu den \textit{Dynamica}\cite{01354} werde von Bodenhausen\indexN{\textls{Bodenhausen} (Bodenus, Bodenausen), Rudolf Christian, Freiherr v. 1698} durch Friedrich von Walter\indexN{\textls{Walter}, Friedrich von gest. 1718} überbracht werden (Brief an R.\,C.\,v.\,Bodenhausen, 20.02.1690,\cite{01527} \textit{LSB} III,\;4 N.\;236, S.\;462.21-463.5, 464.6-7 und S.\;464 Anm. zu N.\;236). Am 4. März bestätigt von Bodenhausen\indexN{\textls{Bodenhausen} (Bodenus, Bodenausen), Rudolf Christian, Freiherr v. 1698}, dieses „Päcklein“ durch von Walter\indexN{\textls{Walter}, Friedrich von gest. 1718} erhalten zu haben, so dass Leibniz davon ausgehen kann, von Bodenhausen\indexN{\textls{Bodenhausen} (Bodenus, Bodenausen), Rudolf Christian, Freiherr v. 1698} habe darin „auch einen bogen gefunden“, der die „ratio Thermometri exacte in gradus dividendi“ enthalte (Leibniz an R.\,C.\,v.\,Bodenhausen, 18.03.1690, \textit{LSB} III,\;4 N.\;245\cite{01392}, S.\;485.1-11; R.\,C.\,v.\,Bodenhausen an Leibniz, 04.03.1690, \textit{LSB} III,\;4 N.\;241\cite{01389}, S.\;470.13-16). Im Folgenden zählt von Bodenhausen\indexN{\textls{Bodenhausen} (Bodenus, Bodenausen), Rudolf Christian, Freiherr v. 1698}  die „ratio dividendi Thermometrum“ dennoch zu denjenigen Nachträgen, die Leibniz zwar angekündigt habe, die ihm aber noch fehlten (R.\,C.\,v.\,Bodenhausen an Leibniz, 29.07.1690, \textit{LSB} III,\;4 N.\;268\cite{01393}, S.\;539.5-7). Leibniz bekräftigt hingegen, dass von Bodenhausen\indexN{\textls{Bodenhausen} (Bodenus, Bodenausen), Rudolf Christian, Freiherr v. 1698} dieses Stück schon bekommen haben müsse  (Leibniz an R.\,C.\,v.\,Bodenhausen, Aug. 1690,\cite{01528} \textit{LSB} III,\;4 N.\;272, S.\;553.10-11), 
woraufhin von Bodenhausen\indexN{\textls{Bodenhausen} (Bodenus, Bodenausen), Rudolf Christian, Freiherr v. 1698} ihm erklärt, dass es in der durch von Walter\indexN{\textls{Walter}, Friedrich von gest. 1718} zugestellten Sendung nicht enthalten gewesen sei (R.\,C.\,v.\,Bodenhausen an Leibniz, 16.09.1690,\cite{01535} \textit{LSB} III,\;4 N.\;273, S.\;558.2-6). 
Auch wenn er sich einen Verlust nicht erklären könne, hält Leibniz auf diese Nachricht hin das „problema Thermometri aequaliter dividendi“ für verloren und für schwer zu ersetzen (Leibniz an R.\,C.\,v.\,Bodenhausen, 26.10./05.11.1690, \textit{LSB} III,\;4 N.\;285\cite{02132}, S.\;630.5-8). Noch in fünf weiteren Briefen erinnert von Bodenhausen\indexN{\textls{Bodenhausen} (Bodenus, Bodenausen), Rudolf Christian, Freiherr v. 1698} daran, dass ihm die „constructio Thermometri“ immer noch fehle (R.\,C.\,v.\,Bodenhausen an Leibniz:  \textit{LSB} III,\;4 N.\;288, S.\;649.14-16 (11.11.1690)\cite{01530}; \textit{LSB} III,\;4 N.\;295, S.\;681.2-4 (12.12.1690)\cite{01531}; \textit{LSB} III,\;5 N.\;3, S.\;31.1-8 (19.01.1691)\cite{01532}; \textit{LSB} III,\;5 N.\;25, S.\;126.8 (23.06.1691)\cite{01394}; \textit{LSB} III,\;5 N.\;125, S.\;474.22-23 (16.01.1693)\cite{01533}). Leibniz gesteht ihm 1693 ein, nicht zu wissen wo seine „Meditationes de Thermometro accurate dividendo“ stecken könnten, hegt aber die Hoffnung, dass sie ihm unter seinen Papieren dereinst wieder „unter die hand kommen müßen“; zugleich macht er von Bodenhausen\indexN{\textls{Bodenhausen} (Bodenus, Bodenausen), Rudolf Christian, Freiherr v. 1698} Mut, solch eine Gradeinteilung des Thermometers selbst finden zu können, da dies nicht schwer sei (Leibniz an R.\,C.\,v.\,Bodenhausen, 12./22.07.1693,\cite{01529} \textit{LSB} III,\;5 N.\;171, S.\;601.4-10). Offenbar unternimmt von Bodenhausen hierauf einen eigenständigen Versuch dazu, bei dem er Leibniz um Hilfe bittet (R.\,C.\,v.\,Bodenhausen an Leibniz, 05.09.1693,\cite{01534} \textit{LSB} III,\;5 N.\;182, S.\;626.1-6).
\pend
\pstart
Es ist anzunehmen, dass es sich bei N.\;\ref{60270} um jenes Stück für die \textit{Dynamica}\cite{01354} handelt, das von Bodenhausen\indexN{\textls{Bodenhausen} (Bodenus, Bodenausen), Rudolf Christian, Freiherr v. 1698} nie erhalten hat. Damit war es, anders als Leibniz und von Bodenhausen\indexN{\textls{Bodenhausen} (Bodenus, Bodenausen), Rudolf Christian, Freiherr v. 1698} in den Jahren 1690 bis 1693 glaubten, nicht verloren gegangen. Wenn es sich bei N.\;\ref{60270}  um dasjenige Stück handelt, das Leibniz seiner Überzeugung nach an von Bodenhausen\indexN{\textls{Bodenhausen} (Bodenus, Bodenausen), Rudolf Christian, Freiherr v. 1698} geschickt hatte, muss er seine Arbeit daran spätestens zu dem Zeitpunkt beendet haben, als er davon ausging, es an von Walter\indexN{\textls{Walter}, Friedrich von gest. 1718} übergeben zu haben: Als \textit{terminus ante quem} ergibt sich daraus für N.\;\ref{60270} der 20. Februar 1690.
Nicht auszuschließen ist, dass Leibniz später selbst das womöglich nur vermeintlich abgeschickte Stück unter seinen Papieren (wie er erhoffte) wiederentdeckte; für eine weitere Bearbeitung desselben, die dann nach dem 22. Juli 1693 erfolgt sein müsste, oder für spätere Hinweise auf dieses Stück gibt es bislang aber keine Anhaltspunkte.
\pend
\pstart
N.\;\ref{60326} ist allem Anschein nach ebenfalls auf römischem Papier geschrieben und dürfte in demselben Zeitraum entstanden sein wie \ref{60270}. N.\;\ref{60326} diente Leibniz als Arbeitsgrundlage für seine Ausarbeitung des Themas in \ref{60270}, die in mehreren Anläufen erfolgt. 
Leibniz führt in N.\;\ref{60326} Rechnungen aus, deren Ergebnisse er für N.\;\ref{60270} übernimmt (\refpassage{LH_37_04_169-170_20a}{LH_37_04_169-170_20b}), beruft sich in N.\;\ref{60270} auf seine Ausführungen in N.\;\ref{60326} (\refpassage{LH_37_04_169-170_19a}{LH_37_04_169-170_19b}); in N.\;\ref{60270} arbeitet er anfangs noch mit einer frühen Fassung von Fig.~2, die sich nur in N.\;\ref{60326} erkennen lässt (\refnote{LH_37_04_169-170_21ab}), bevor er diese Zeichnung korrigiert und eine neue dafür in N.\;\ref{60270} (Fig.~2) erstellt; er benutzt in N.\;\ref{60270} unvermittelt Größen ($\varepsilon$, \textit{V} bzw. \textit{Y}), die er in N.\;\ref{60326} herleitet; die Bezeichnung für die Größen \textit{V} bzw. \textit{Y} ändert Leibniz im Laufe der Ausarbeitung von N.\;\ref{60270} und dokumentiert dies nachträglich in N.\;\ref{60326} (\refpassage{LH_37_05_108_01a}{LH_37_05_108_01b}).%
\pend
\selectlanguage{latin}%
\frenchspacing%
\normalsize%






